% ============================================================
% plano-inclinado.tex
% Macros TikZ para dibujar un plano inclinado y trabajar con
% puntos, tangente, normal, proyecciones y tramo rugoso.
%
% Uso en main.tex:
%   \usetikzlibrary{calc,angles,quotes}
%   % ============================================================
% plano-inclinado.tex
% Macros TikZ para dibujar un plano inclinado y trabajar con
% puntos, tangente, normal, proyecciones y tramo rugoso.
%
% Uso en main.tex:
%   \usetikzlibrary{calc,angles,quotes}
%   % ============================================================
% plano-inclinado.tex
% Macros TikZ para dibujar un plano inclinado y trabajar con
% puntos, tangente, normal, proyecciones y tramo rugoso.
%
% Uso en main.tex:
%   \usetikzlibrary{calc,angles,quotes}
%   % ============================================================
% plano-inclinado.tex
% Macros TikZ para dibujar un plano inclinado y trabajar con
% puntos, tangente, normal, proyecciones y tramo rugoso.
%
% Uso en main.tex:
%   \usetikzlibrary{calc,angles,quotes}
%   \input{plano-inclinado.tex}
%
% Por ahora se administra UN plano por tikzpicture, porque las
% coordenadas internas usan nombres PI-... fijos.
% ============================================================

\makeatletter

% Variables internas para guardar las coordenadas y comparar
\newdimen\PI@xOne \newdimen\PI@yOne
\newdimen\PI@xTwo \newdimen\PI@yTwo
\newdimen\PI@xA   \newdimen\PI@yA
\newdimen\PI@xB   \newdimen\PI@yB

% Estado y valores por defecto de las opciones
\newif\ifPI@showprojection
\newif\ifPI@showfloor
\newif\ifPI@showangle

\pgfkeys{
  /plano/.is family,
  /plano,
  order/.store in=\PI@order,
  show projection/.is if=PI@showprojection,
  show projection/.default=true,
  show floor/.is if=PI@showfloor,
  show floor/.default=true,
  show angle/.is if=PI@showangle,
  show angle/.default=true,
  angle/.store in=\PI@anglelabel,
  angle radius/.store in=\PI@angleradius,
  projection style/.store in=\PI@projectionstyle,
  floor style/.store in=\PI@floorstyle,
  plane style/.store in=\PI@planestyle
}

% \planoInclinado[opciones]{punto1}{punto2}
\newcommand{\planoInclinado}[3][]{%
  % Restablecer todos los valores en cada llamada
  \PI@showprojectionfalse
  \PI@showfloorfalse
  \PI@showanglefalse
  \def\PI@order{bottomtop}%
  \def\PI@anglelabel{\theta}%
  \def\PI@angleradius{7mm}%
  \def\PI@projectionstyle{dashed}%
  \def\PI@floorstyle{dotted}%
  \def\PI@planestyle{thick}%
  \pgfkeys{/plano/.cd,#1}%
  
  % Coordenadas recibidas
  \coordinate (PI-one) at (#2);
  \coordinate (PI-two) at (#3);
  
  % Extraer coordenadas cartesianas para evaluar el sentido
  \pgfextractx{\PI@xOne}{\pgfpointanchor{PI-one}{center}}
  \pgfextracty{\PI@yOne}{\pgfpointanchor{PI-one}{center}}
  \pgfextractx{\PI@xTwo}{\pgfpointanchor{PI-two}{center}}
  \pgfextracty{\PI@yTwo}{\pgfpointanchor{PI-two}{center}}

  \def\PI@ordertopbottom{topbottom}%
  \def\PI@ordertopdown{topdown}%
  
  \ifx\PI@order\PI@ordertopbottom
    \coordinate (PI-A) at (PI-two);
    \coordinate (PI-B) at (PI-one);
    \PI@xA=\PI@xTwo \PI@xB=\PI@xOne
  \else\ifx\PI@order\PI@ordertopdown
    \coordinate (PI-A) at (PI-two);
    \coordinate (PI-B) at (PI-one);
    \PI@xA=\PI@xTwo \PI@xB=\PI@xOne
  \else
    \def\PI@orderbottomup{bottomup}%
    \ifx\PI@order\PI@orderbottomup
      \coordinate (PI-A) at (PI-one);
      \coordinate (PI-B) at (PI-two);
      \PI@xA=\PI@xOne \PI@xB=\PI@xTwo
    \else
      \coordinate (PI-A) at (PI-one);
      \coordinate (PI-B) at (PI-two);
      \PI@xA=\PI@xOne \PI@xB=\PI@xTwo
    \fi
  \fi\fi
  
  % Proyección horizontal de B sobre la altura de A
  \coordinate (PI-P) at (PI-B |- PI-A);
  
  % Geometría principal
  \draw[\PI@planestyle] (PI-A) -- (PI-B);
  \ifPI@showprojection
    \draw[\PI@projectionstyle] (PI-B) -- (PI-P);
  \fi
  \ifPI@showfloor
    \draw[\PI@floorstyle] (PI-A) -- (PI-P);
  \fi
  
  % Ángulo adaptado al sentido de la rampa
  \ifPI@showangle
    \ifdim\PI@xB>\PI@xA
      % La rampa asciende hacia la derecha
      \pic[
        draw,
        angle radius=\PI@angleradius,
        angle eccentricity=1.3,
        "$\PI@anglelabel$" {font=\small}
      ] {angle = PI-P--PI-A--PI-B};
    \else
      % La rampa asciende hacia la izquierda
      \pic[
        draw,
        angle radius=\PI@angleradius,
        angle eccentricity=1.3,
        "$\PI@anglelabel$" {font=\small}
      ] {angle = PI-B--PI-A--PI-P};
    \fi
  \fi
  
  % Guardar ángulo de la tangente A -> B y el de la normal superior.
  \pgfmathanglebetweenpoints{\pgfpointanchor{PI-A}{center}}{\pgfpointanchor{PI-B}{center}}%
  \edef\PI@tangentangle{\pgfmathresult}%
  \ifdim\PI@tangentangle pt>90pt
    \pgfmathsetmacro{\PI@normalangle}{\PI@tangentangle-90}%
  \else
    \pgfmathsetmacro{\PI@normalangle}{\PI@tangentangle+90}%
  \fi
  
  % Componentes unitarias disponibles para cálculos posteriores.
  \pgfmathsetmacro{\PI@tx}{cos(\PI@tangentangle)}%
  \pgfmathsetmacro{\PI@ty}{sin(\PI@tangentangle)}%
  \pgfmathsetmacro{\PI@nx}{cos(\PI@normalangle)}%
  \pgfmathsetmacro{\PI@ny}{sin(\PI@normalangle)}%
  \coordinate (PI-tangente) at (\PI@tx,\PI@ty);
  \coordinate (PI-normal) at (\PI@nx,\PI@ny);
}

% \puntoRampa{t}{nombre}
\newcommand{\puntoRampa}[2]{%
  \coordinate (#2) at ($(PI-A)!#1!(PI-B)$);%
}

% Atajo solicitado: calcula el punto y lo guarda como (PI-punto).
\newcommand{\atRampa}[1]{\puntoRampa{#1}{PI-punto}}

% \tangenteRampa{nombre} y \normalRampa{nombre}
\newcommand{\tangenteRampa}[1]{\coordinate (#1) at (\PI@tx,\PI@ty);}
\newcommand{\normalRampa}[1]{\coordinate (#1) at (\PI@nx,\PI@ny);}

% \proyeccionPiso{coordenada-origen}{nombre-destino}
\newcommand{\proyeccionPiso}[2]{\coordinate (#2) at (#1 |- PI-A);}

% \centroCirculoRampa{t}{radio}{nombre-centro}
\newcommand{\centroCirculoRampa}[3]{%
  \puntoRampa{#1}{PI-contacto-temp}%
  \coordinate (#3) at
    ($(PI-contacto-temp)+({(#2)*\PI@nx},{(#2)*\PI@ny})$);%
}

% \circuloRampa[estilo de dibujo]{t}{radio}{nombre-centro}
\newcommand{\circuloRampa}[4][thick]{%
    \centroCirculoRampa{#2}{#3}{#4}%
    \draw[#1] (#4) circle (#3);%
}

% \tramoRugoso[estilo]{t-inicial}{t-final}
\newcommand{\tramoRugoso}[3][very thick]{%
  \puntoRampa{#2}{PI-rugoso-inicio}%
  \puntoRampa{#3}{PI-rugoso-fin}%
  \draw[#1] (PI-rugoso-inicio) -- (PI-rugoso-fin);%
}

% \flechaTangente[estilo]{punto}{longitud-firmada}
\newcommand{\flechaTangente}[3][->,thick]{%
  \draw[#1] (#2) -- ++({(#3)*\PI@tx},{(#3)*\PI@ty});%
}

% \flechaNormal[estilo]{punto}{longitud-firmada}
\newcommand{\flechaNormal}[3][->,thick]{%
  \draw[#1] (#2) -- ++({(#3)*\PI@nx},{(#3)*\PI@ny});%
}

\makeatother

%
% Por ahora se administra UN plano por tikzpicture, porque las
% coordenadas internas usan nombres PI-... fijos.
% ============================================================

\makeatletter

% Variables internas para guardar las coordenadas y comparar
\newdimen\PI@xOne \newdimen\PI@yOne
\newdimen\PI@xTwo \newdimen\PI@yTwo
\newdimen\PI@xA   \newdimen\PI@yA
\newdimen\PI@xB   \newdimen\PI@yB

% Estado y valores por defecto de las opciones
\newif\ifPI@showprojection
\newif\ifPI@showfloor
\newif\ifPI@showangle

\pgfkeys{
  /plano/.is family,
  /plano,
  order/.store in=\PI@order,
  show projection/.is if=PI@showprojection,
  show projection/.default=true,
  show floor/.is if=PI@showfloor,
  show floor/.default=true,
  show angle/.is if=PI@showangle,
  show angle/.default=true,
  angle/.store in=\PI@anglelabel,
  angle radius/.store in=\PI@angleradius,
  projection style/.store in=\PI@projectionstyle,
  floor style/.store in=\PI@floorstyle,
  plane style/.store in=\PI@planestyle
}

% \planoInclinado[opciones]{punto1}{punto2}
\newcommand{\planoInclinado}[3][]{%
  % Restablecer todos los valores en cada llamada
  \PI@showprojectionfalse
  \PI@showfloorfalse
  \PI@showanglefalse
  \def\PI@order{bottomtop}%
  \def\PI@anglelabel{\theta}%
  \def\PI@angleradius{7mm}%
  \def\PI@projectionstyle{dashed}%
  \def\PI@floorstyle{dotted}%
  \def\PI@planestyle{thick}%
  \pgfkeys{/plano/.cd,#1}%
  
  % Coordenadas recibidas
  \coordinate (PI-one) at (#2);
  \coordinate (PI-two) at (#3);
  
  % Extraer coordenadas cartesianas para evaluar el sentido
  \pgfextractx{\PI@xOne}{\pgfpointanchor{PI-one}{center}}
  \pgfextracty{\PI@yOne}{\pgfpointanchor{PI-one}{center}}
  \pgfextractx{\PI@xTwo}{\pgfpointanchor{PI-two}{center}}
  \pgfextracty{\PI@yTwo}{\pgfpointanchor{PI-two}{center}}

  \def\PI@ordertopbottom{topbottom}%
  \def\PI@ordertopdown{topdown}%
  
  \ifx\PI@order\PI@ordertopbottom
    \coordinate (PI-A) at (PI-two);
    \coordinate (PI-B) at (PI-one);
    \PI@xA=\PI@xTwo \PI@xB=\PI@xOne
  \else\ifx\PI@order\PI@ordertopdown
    \coordinate (PI-A) at (PI-two);
    \coordinate (PI-B) at (PI-one);
    \PI@xA=\PI@xTwo \PI@xB=\PI@xOne
  \else
    \def\PI@orderbottomup{bottomup}%
    \ifx\PI@order\PI@orderbottomup
      \coordinate (PI-A) at (PI-one);
      \coordinate (PI-B) at (PI-two);
      \PI@xA=\PI@xOne \PI@xB=\PI@xTwo
    \else
      \coordinate (PI-A) at (PI-one);
      \coordinate (PI-B) at (PI-two);
      \PI@xA=\PI@xOne \PI@xB=\PI@xTwo
    \fi
  \fi\fi
  
  % Proyección horizontal de B sobre la altura de A
  \coordinate (PI-P) at (PI-B |- PI-A);
  
  % Geometría principal
  \draw[\PI@planestyle] (PI-A) -- (PI-B);
  \ifPI@showprojection
    \draw[\PI@projectionstyle] (PI-B) -- (PI-P);
  \fi
  \ifPI@showfloor
    \draw[\PI@floorstyle] (PI-A) -- (PI-P);
  \fi
  
  % Ángulo adaptado al sentido de la rampa
  \ifPI@showangle
    \ifdim\PI@xB>\PI@xA
      % La rampa asciende hacia la derecha
      \pic[
        draw,
        angle radius=\PI@angleradius,
        angle eccentricity=1.3,
        "$\PI@anglelabel$" {font=\small}
      ] {angle = PI-P--PI-A--PI-B};
    \else
      % La rampa asciende hacia la izquierda
      \pic[
        draw,
        angle radius=\PI@angleradius,
        angle eccentricity=1.3,
        "$\PI@anglelabel$" {font=\small}
      ] {angle = PI-B--PI-A--PI-P};
    \fi
  \fi
  
  % Guardar ángulo de la tangente A -> B y el de la normal superior.
  \pgfmathanglebetweenpoints{\pgfpointanchor{PI-A}{center}}{\pgfpointanchor{PI-B}{center}}%
  \edef\PI@tangentangle{\pgfmathresult}%
  \ifdim\PI@tangentangle pt>90pt
    \pgfmathsetmacro{\PI@normalangle}{\PI@tangentangle-90}%
  \else
    \pgfmathsetmacro{\PI@normalangle}{\PI@tangentangle+90}%
  \fi
  
  % Componentes unitarias disponibles para cálculos posteriores.
  \pgfmathsetmacro{\PI@tx}{cos(\PI@tangentangle)}%
  \pgfmathsetmacro{\PI@ty}{sin(\PI@tangentangle)}%
  \pgfmathsetmacro{\PI@nx}{cos(\PI@normalangle)}%
  \pgfmathsetmacro{\PI@ny}{sin(\PI@normalangle)}%
  \coordinate (PI-tangente) at (\PI@tx,\PI@ty);
  \coordinate (PI-normal) at (\PI@nx,\PI@ny);
}

% \puntoRampa{t}{nombre}
\newcommand{\puntoRampa}[2]{%
  \coordinate (#2) at ($(PI-A)!#1!(PI-B)$);%
}

% Atajo solicitado: calcula el punto y lo guarda como (PI-punto).
\newcommand{\atRampa}[1]{\puntoRampa{#1}{PI-punto}}

% \tangenteRampa{nombre} y \normalRampa{nombre}
\newcommand{\tangenteRampa}[1]{\coordinate (#1) at (\PI@tx,\PI@ty);}
\newcommand{\normalRampa}[1]{\coordinate (#1) at (\PI@nx,\PI@ny);}

% \proyeccionPiso{coordenada-origen}{nombre-destino}
\newcommand{\proyeccionPiso}[2]{\coordinate (#2) at (#1 |- PI-A);}

% \centroCirculoRampa{t}{radio}{nombre-centro}
\newcommand{\centroCirculoRampa}[3]{%
  \puntoRampa{#1}{PI-contacto-temp}%
  \coordinate (#3) at
    ($(PI-contacto-temp)+({(#2)*\PI@nx},{(#2)*\PI@ny})$);%
}

% \circuloRampa[estilo de dibujo]{t}{radio}{nombre-centro}
\newcommand{\circuloRampa}[4][thick]{%
    \centroCirculoRampa{#2}{#3}{#4}%
    \draw[#1] (#4) circle (#3);%
}

% \tramoRugoso[estilo]{t-inicial}{t-final}
\newcommand{\tramoRugoso}[3][very thick]{%
  \puntoRampa{#2}{PI-rugoso-inicio}%
  \puntoRampa{#3}{PI-rugoso-fin}%
  \draw[#1] (PI-rugoso-inicio) -- (PI-rugoso-fin);%
}

% \flechaTangente[estilo]{punto}{longitud-firmada}
\newcommand{\flechaTangente}[3][->,thick]{%
  \draw[#1] (#2) -- ++({(#3)*\PI@tx},{(#3)*\PI@ty});%
}

% \flechaNormal[estilo]{punto}{longitud-firmada}
\newcommand{\flechaNormal}[3][->,thick]{%
  \draw[#1] (#2) -- ++({(#3)*\PI@nx},{(#3)*\PI@ny});%
}

\makeatother

%
% Por ahora se administra UN plano por tikzpicture, porque las
% coordenadas internas usan nombres PI-... fijos.
% ============================================================

\makeatletter

% Variables internas para guardar las coordenadas y comparar
\newdimen\PI@xOne \newdimen\PI@yOne
\newdimen\PI@xTwo \newdimen\PI@yTwo
\newdimen\PI@xA   \newdimen\PI@yA
\newdimen\PI@xB   \newdimen\PI@yB

% Estado y valores por defecto de las opciones
\newif\ifPI@showprojection
\newif\ifPI@showfloor
\newif\ifPI@showangle

\pgfkeys{
  /plano/.is family,
  /plano,
  order/.store in=\PI@order,
  show projection/.is if=PI@showprojection,
  show projection/.default=true,
  show floor/.is if=PI@showfloor,
  show floor/.default=true,
  show angle/.is if=PI@showangle,
  show angle/.default=true,
  angle/.store in=\PI@anglelabel,
  angle radius/.store in=\PI@angleradius,
  projection style/.store in=\PI@projectionstyle,
  floor style/.store in=\PI@floorstyle,
  plane style/.store in=\PI@planestyle
}

% \planoInclinado[opciones]{punto1}{punto2}
\newcommand{\planoInclinado}[3][]{%
  % Restablecer todos los valores en cada llamada
  \PI@showprojectionfalse
  \PI@showfloorfalse
  \PI@showanglefalse
  \def\PI@order{bottomtop}%
  \def\PI@anglelabel{\theta}%
  \def\PI@angleradius{7mm}%
  \def\PI@projectionstyle{dashed}%
  \def\PI@floorstyle{dotted}%
  \def\PI@planestyle{thick}%
  \pgfkeys{/plano/.cd,#1}%
  
  % Coordenadas recibidas
  \coordinate (PI-one) at (#2);
  \coordinate (PI-two) at (#3);
  
  % Extraer coordenadas cartesianas para evaluar el sentido
  \pgfextractx{\PI@xOne}{\pgfpointanchor{PI-one}{center}}
  \pgfextracty{\PI@yOne}{\pgfpointanchor{PI-one}{center}}
  \pgfextractx{\PI@xTwo}{\pgfpointanchor{PI-two}{center}}
  \pgfextracty{\PI@yTwo}{\pgfpointanchor{PI-two}{center}}

  \def\PI@ordertopbottom{topbottom}%
  \def\PI@ordertopdown{topdown}%
  
  \ifx\PI@order\PI@ordertopbottom
    \coordinate (PI-A) at (PI-two);
    \coordinate (PI-B) at (PI-one);
    \PI@xA=\PI@xTwo \PI@xB=\PI@xOne
  \else\ifx\PI@order\PI@ordertopdown
    \coordinate (PI-A) at (PI-two);
    \coordinate (PI-B) at (PI-one);
    \PI@xA=\PI@xTwo \PI@xB=\PI@xOne
  \else
    \def\PI@orderbottomup{bottomup}%
    \ifx\PI@order\PI@orderbottomup
      \coordinate (PI-A) at (PI-one);
      \coordinate (PI-B) at (PI-two);
      \PI@xA=\PI@xOne \PI@xB=\PI@xTwo
    \else
      \coordinate (PI-A) at (PI-one);
      \coordinate (PI-B) at (PI-two);
      \PI@xA=\PI@xOne \PI@xB=\PI@xTwo
    \fi
  \fi\fi
  
  % Proyección horizontal de B sobre la altura de A
  \coordinate (PI-P) at (PI-B |- PI-A);
  
  % Geometría principal
  \draw[\PI@planestyle] (PI-A) -- (PI-B);
  \ifPI@showprojection
    \draw[\PI@projectionstyle] (PI-B) -- (PI-P);
  \fi
  \ifPI@showfloor
    \draw[\PI@floorstyle] (PI-A) -- (PI-P);
  \fi
  
  % Ángulo adaptado al sentido de la rampa
  \ifPI@showangle
    \ifdim\PI@xB>\PI@xA
      % La rampa asciende hacia la derecha
      \pic[
        draw,
        angle radius=\PI@angleradius,
        angle eccentricity=1.3,
        "$\PI@anglelabel$" {font=\small}
      ] {angle = PI-P--PI-A--PI-B};
    \else
      % La rampa asciende hacia la izquierda
      \pic[
        draw,
        angle radius=\PI@angleradius,
        angle eccentricity=1.3,
        "$\PI@anglelabel$" {font=\small}
      ] {angle = PI-B--PI-A--PI-P};
    \fi
  \fi
  
  % Guardar ángulo de la tangente A -> B y el de la normal superior.
  \pgfmathanglebetweenpoints{\pgfpointanchor{PI-A}{center}}{\pgfpointanchor{PI-B}{center}}%
  \edef\PI@tangentangle{\pgfmathresult}%
  \ifdim\PI@tangentangle pt>90pt
    \pgfmathsetmacro{\PI@normalangle}{\PI@tangentangle-90}%
  \else
    \pgfmathsetmacro{\PI@normalangle}{\PI@tangentangle+90}%
  \fi
  
  % Componentes unitarias disponibles para cálculos posteriores.
  \pgfmathsetmacro{\PI@tx}{cos(\PI@tangentangle)}%
  \pgfmathsetmacro{\PI@ty}{sin(\PI@tangentangle)}%
  \pgfmathsetmacro{\PI@nx}{cos(\PI@normalangle)}%
  \pgfmathsetmacro{\PI@ny}{sin(\PI@normalangle)}%
  \coordinate (PI-tangente) at (\PI@tx,\PI@ty);
  \coordinate (PI-normal) at (\PI@nx,\PI@ny);
}

% \puntoRampa{t}{nombre}
\newcommand{\puntoRampa}[2]{%
  \coordinate (#2) at ($(PI-A)!#1!(PI-B)$);%
}

% Atajo solicitado: calcula el punto y lo guarda como (PI-punto).
\newcommand{\atRampa}[1]{\puntoRampa{#1}{PI-punto}}

% \tangenteRampa{nombre} y \normalRampa{nombre}
\newcommand{\tangenteRampa}[1]{\coordinate (#1) at (\PI@tx,\PI@ty);}
\newcommand{\normalRampa}[1]{\coordinate (#1) at (\PI@nx,\PI@ny);}

% \proyeccionPiso{coordenada-origen}{nombre-destino}
\newcommand{\proyeccionPiso}[2]{\coordinate (#2) at (#1 |- PI-A);}

% \centroCirculoRampa{t}{radio}{nombre-centro}
\newcommand{\centroCirculoRampa}[3]{%
  \puntoRampa{#1}{PI-contacto-temp}%
  \coordinate (#3) at
    ($(PI-contacto-temp)+({(#2)*\PI@nx},{(#2)*\PI@ny})$);%
}

% \circuloRampa[estilo de dibujo]{t}{radio}{nombre-centro}
\newcommand{\circuloRampa}[4][thick]{%
    \centroCirculoRampa{#2}{#3}{#4}%
    \draw[#1] (#4) circle (#3);%
}

% \tramoRugoso[estilo]{t-inicial}{t-final}
\newcommand{\tramoRugoso}[3][very thick]{%
  \puntoRampa{#2}{PI-rugoso-inicio}%
  \puntoRampa{#3}{PI-rugoso-fin}%
  \draw[#1] (PI-rugoso-inicio) -- (PI-rugoso-fin);%
}

% \flechaTangente[estilo]{punto}{longitud-firmada}
\newcommand{\flechaTangente}[3][->,thick]{%
  \draw[#1] (#2) -- ++({(#3)*\PI@tx},{(#3)*\PI@ty});%
}

% \flechaNormal[estilo]{punto}{longitud-firmada}
\newcommand{\flechaNormal}[3][->,thick]{%
  \draw[#1] (#2) -- ++({(#3)*\PI@nx},{(#3)*\PI@ny});%
}

\makeatother

%
% Por ahora se administra UN plano por tikzpicture, porque las
% coordenadas internas usan nombres PI-... fijos.
% ============================================================

\makeatletter

% Variables internas para guardar las coordenadas y comparar
\newdimen\PI@xOne \newdimen\PI@yOne
\newdimen\PI@xTwo \newdimen\PI@yTwo
\newdimen\PI@xA   \newdimen\PI@yA
\newdimen\PI@xB   \newdimen\PI@yB

% Estado y valores por defecto de las opciones
\newif\ifPI@showprojection
\newif\ifPI@showfloor
\newif\ifPI@showangle

\pgfkeys{
  /plano/.is family,
  /plano,
  order/.store in=\PI@order,
  show projection/.is if=PI@showprojection,
  show projection/.default=true,
  show floor/.is if=PI@showfloor,
  show floor/.default=true,
  show angle/.is if=PI@showangle,
  show angle/.default=true,
  angle/.store in=\PI@anglelabel,
  angle radius/.store in=\PI@angleradius,
  projection style/.store in=\PI@projectionstyle,
  floor style/.store in=\PI@floorstyle,
  plane style/.store in=\PI@planestyle
}

% \planoInclinado[opciones]{punto1}{punto2}
\newcommand{\planoInclinado}[3][]{%
  % Restablecer todos los valores en cada llamada
  \PI@showprojectionfalse
  \PI@showfloorfalse
  \PI@showanglefalse
  \def\PI@order{bottomtop}%
  \def\PI@anglelabel{\theta}%
  \def\PI@angleradius{7mm}%
  \def\PI@projectionstyle{dashed}%
  \def\PI@floorstyle{dotted}%
  \def\PI@planestyle{thick}%
  \pgfkeys{/plano/.cd,#1}%
  
  % Coordenadas recibidas
  \coordinate (PI-one) at (#2);
  \coordinate (PI-two) at (#3);
  
  % Extraer coordenadas cartesianas para evaluar el sentido
  \pgfextractx{\PI@xOne}{\pgfpointanchor{PI-one}{center}}
  \pgfextracty{\PI@yOne}{\pgfpointanchor{PI-one}{center}}
  \pgfextractx{\PI@xTwo}{\pgfpointanchor{PI-two}{center}}
  \pgfextracty{\PI@yTwo}{\pgfpointanchor{PI-two}{center}}

  \def\PI@ordertopbottom{topbottom}%
  \def\PI@ordertopdown{topdown}%
  
  \ifx\PI@order\PI@ordertopbottom
    \coordinate (PI-A) at (PI-two);
    \coordinate (PI-B) at (PI-one);
    \PI@xA=\PI@xTwo \PI@xB=\PI@xOne
  \else\ifx\PI@order\PI@ordertopdown
    \coordinate (PI-A) at (PI-two);
    \coordinate (PI-B) at (PI-one);
    \PI@xA=\PI@xTwo \PI@xB=\PI@xOne
  \else
    \def\PI@orderbottomup{bottomup}%
    \ifx\PI@order\PI@orderbottomup
      \coordinate (PI-A) at (PI-one);
      \coordinate (PI-B) at (PI-two);
      \PI@xA=\PI@xOne \PI@xB=\PI@xTwo
    \else
      \coordinate (PI-A) at (PI-one);
      \coordinate (PI-B) at (PI-two);
      \PI@xA=\PI@xOne \PI@xB=\PI@xTwo
    \fi
  \fi\fi
  
  % Proyección horizontal de B sobre la altura de A
  \coordinate (PI-P) at (PI-B |- PI-A);
  
  % Geometría principal
  \draw[\PI@planestyle] (PI-A) -- (PI-B);
  \ifPI@showprojection
    \draw[\PI@projectionstyle] (PI-B) -- (PI-P);
  \fi
  \ifPI@showfloor
    \draw[\PI@floorstyle] (PI-A) -- (PI-P);
  \fi
  
  % Ángulo adaptado al sentido de la rampa
  \ifPI@showangle
    \ifdim\PI@xB>\PI@xA
      % La rampa asciende hacia la derecha
      \pic[
        draw,
        angle radius=\PI@angleradius,
        angle eccentricity=1.3,
        "$\PI@anglelabel$" {font=\small}
      ] {angle = PI-P--PI-A--PI-B};
    \else
      % La rampa asciende hacia la izquierda
      \pic[
        draw,
        angle radius=\PI@angleradius,
        angle eccentricity=1.3,
        "$\PI@anglelabel$" {font=\small}
      ] {angle = PI-B--PI-A--PI-P};
    \fi
  \fi
  
  % Guardar ángulo de la tangente A -> B y el de la normal superior.
  \pgfmathanglebetweenpoints{\pgfpointanchor{PI-A}{center}}{\pgfpointanchor{PI-B}{center}}%
  \edef\PI@tangentangle{\pgfmathresult}%
  \ifdim\PI@tangentangle pt>90pt
    \pgfmathsetmacro{\PI@normalangle}{\PI@tangentangle-90}%
  \else
    \pgfmathsetmacro{\PI@normalangle}{\PI@tangentangle+90}%
  \fi
  
  % Componentes unitarias disponibles para cálculos posteriores.
  \pgfmathsetmacro{\PI@tx}{cos(\PI@tangentangle)}%
  \pgfmathsetmacro{\PI@ty}{sin(\PI@tangentangle)}%
  \pgfmathsetmacro{\PI@nx}{cos(\PI@normalangle)}%
  \pgfmathsetmacro{\PI@ny}{sin(\PI@normalangle)}%
  \coordinate (PI-tangente) at (\PI@tx,\PI@ty);
  \coordinate (PI-normal) at (\PI@nx,\PI@ny);
}

% \puntoRampa{t}{nombre}
\newcommand{\puntoRampa}[2]{%
  \coordinate (#2) at ($(PI-A)!#1!(PI-B)$);%
}

% Atajo solicitado: calcula el punto y lo guarda como (PI-punto).
\newcommand{\atRampa}[1]{\puntoRampa{#1}{PI-punto}}

% \tangenteRampa{nombre} y \normalRampa{nombre}
\newcommand{\tangenteRampa}[1]{\coordinate (#1) at (\PI@tx,\PI@ty);}
\newcommand{\normalRampa}[1]{\coordinate (#1) at (\PI@nx,\PI@ny);}

% \proyeccionPiso{coordenada-origen}{nombre-destino}
\newcommand{\proyeccionPiso}[2]{\coordinate (#2) at (#1 |- PI-A);}

% \centroCirculoRampa{t}{radio}{nombre-centro}
\newcommand{\centroCirculoRampa}[3]{%
  \puntoRampa{#1}{PI-contacto-temp}%
  \coordinate (#3) at
    ($(PI-contacto-temp)+({(#2)*\PI@nx},{(#2)*\PI@ny})$);%
}

% \circuloRampa[estilo de dibujo]{t}{radio}{nombre-centro}
\newcommand{\circuloRampa}[4][thick]{%
    \centroCirculoRampa{#2}{#3}{#4}%
    \draw[#1] (#4) circle (#3);%
}

% \tramoRugoso[estilo]{t-inicial}{t-final}
\newcommand{\tramoRugoso}[3][very thick]{%
  \puntoRampa{#2}{PI-rugoso-inicio}%
  \puntoRampa{#3}{PI-rugoso-fin}%
  \draw[#1] (PI-rugoso-inicio) -- (PI-rugoso-fin);%
}

% \flechaTangente[estilo]{punto}{longitud-firmada}
\newcommand{\flechaTangente}[3][->,thick]{%
  \draw[#1] (#2) -- ++({(#3)*\PI@tx},{(#3)*\PI@ty});%
}

% \flechaNormal[estilo]{punto}{longitud-firmada}
\newcommand{\flechaNormal}[3][->,thick]{%
  \draw[#1] (#2) -- ++({(#3)*\PI@nx},{(#3)*\PI@ny});%
}

\makeatother
